% Source-aligned English for printed p.185.
\clearpage
\SourcePage{185}

\LayerTarget{tohoku:chapter:IV}
\begin{center}
{\large\bfseries Chapter IV\quad Ext of sheaves of modules.}
\end{center}

\LayerTarget{tohoku:section:4.1}
\SectionHeading{4.1. The functors $\mathrm{Hom}_O(A,B)$ and
$\mathbf{Hom}_O(A,B)$.}

\LayerTarget{tohoku:definition:hom-o-and-local-hom}
Let $X$ be a topological space equipped with a sheaf $O$ of rings with
identity, and let $C^O$ be the abelian category of (left) $O$-modules on $X$
(cf. 3.1). If $A$ and $B$ are two $O$-modules, we denote by
$\mathrm{Hom}_O(A,B)$ the group of $O$-homomorphisms from the first to the
second; this group is in fact a module over $\Gamma(O^\natural)$. If $U$
is an arbitrary subset of $X$, put
\[
\mathrm{Hom}_O(U;A,B)=\mathrm{Hom}_{O|U}(A|U,B|U),
\]
where $F|U$ denotes, as usual, the restriction of a sheaf $F$ to $U$.

\LayerTarget{tohoku:definition:sheaf-hom-o}
Of course, if $V\subset U$, there is a natural homomorphism
$\mathrm{Hom}_O(U;A,B)\longrightarrow\mathrm{Hom}_O(V;A,B)$. One immediately
verifies that, when restricted to the open subsets $U$ of $X$, the groups
$\mathrm{Hom}_O(U;A,B)$ form a \textit{sheaf} on $X$, denoted
$\mathbf{Hom}_O(A,B)$. Thus, by definition,

\LayerTarget{tohoku:eq:4.1.1}
\[
\tag{4.1.1}
\Gamma(U,\mathbf{Hom}_O(A,B))=\mathrm{Hom}_O(U;A,B),
\]
and in particular, on taking $U=X$,

\LayerTarget{tohoku:eq:4.1.2}
\[
\tag{4.1.2}
\mathrm{Hom}_O(A,B)=\Gamma(\mathbf{Hom}_O(A,B)).
\]

\LayerTarget{tohoku:remark:sheaf-hom-o:center-module}
Moreover, $\mathbf{Hom}_O(A,B)$ can also be regarded as a sheaf of modules
over the center $O^\natural$ of $O$.

\LayerTarget{tohoku:remark:hom-o:left-exactness}
Recall that $\mathrm{Hom}_O(A,B)$ is an \textit{additive left exact functor}
in the two arguments $A,B\in C^O$, with values in abelian groups. It follows
that $\mathbf{Hom}_O(A,B)$ can also be regarded as a functor left exact in the
two arguments $A,B\in C^O$, with values in the category $C^X$ of abelian
sheaves on $X$, or even with values in the category $C^{O^\natural}$. As
usual, all homomorphisms we write will be ``functorial.'' This section gives
some auxiliary properties of the functors under consideration, preliminary
to the study of their derived functors.

Let $A,B\in C^O$, and let $x\in X$. For every open $U$ containing $x$, there
is an evident homomorphism
$\mathrm{Hom}_O(U;A,B)\longrightarrow\mathrm{Hom}_{O(x)}(A(x),B(x))$; hence,
on passing to the inductive limit over the open neighborhoods of $x$, there is
a natural homomorphism

\LayerTarget{tohoku:eq:4.1.3}
\[
\tag{4.1.3}
\mathbf{Hom}_O(A,B)(x)\longrightarrow\mathrm{Hom}_{O(x)}(A(x),B(x)),
\]
which we shall study.

\LayerTarget{tohoku:definition:finite-type-at-x}
We say that $A$ is \textit{of finite type at $x$} if there is an open
neighborhood $U$ of $x$ such that $A|U$ is isomorphic to a quotient of
$(O|U)^n$, where $n$ is a finite positive integer.

\LayerTarget{tohoku:definition:pseudo-coherent-at-x}
The object $A$ is called \textit{pseudo-coherent at $x$} if there is an open
neighborhood $U$ of $x$ and, \textit{on $U$}, an exact sequence of
homomorphisms $O^m\longrightarrow O^n\longrightarrow A\longrightarrow0$,
where $m,n$ are finite positive integers. We say that $A$ is of finite type
(respectively, pseudo-coherent) if it is so at every point. With this
understood:

\LayerTarget{tohoku:proposition:4.1.1}
\noindent\textsc{Proposition 4.1.1.} \textit{The homomorphism (4.1.3) is a
monomorphism if $A$ is of finite type at $x$, and an isomorphism if $A$ is
pseudo-coherent at $x$.}

\LayerTarget{tohoku:proof:proposition-4.1.1:p1}
Suppose that $A$ is of finite type at $x$. After restricting, if necessary,
to a suitable neighborhood of $x$, there is an exact sequence
$O^n\longrightarrow A\longrightarrow0$, hence an exact sequence
$0\longrightarrow\mathbf{Hom}_O(A,B)\longrightarrow
\mathbf{Hom}_O(O^n,B)$, and an exact sequence
$O^n(x)\longrightarrow A(x)\longrightarrow0$, hence an exact sequence
$0\longrightarrow\mathrm{Hom}_{O(x)}(A(x),B(x))\longrightarrow
\mathrm{Hom}_{O(x)}(O^n(x),B(x))$. From these we obtain a homomorphism of exact
sequences:

\TohokuContinuation{186}
